%% parts/p07/arithmetic-proof-theory.tex --- Phase-4 P7 chapter 7.5
%% Source migration: v1.1 Part VIII, lines 12133--12216.
%% Depth/new-content: D-P8-005, D-P8-006, N-a-049--N-a-052.
%% N-a-049 remains explicitly stubbed at its unmet full-assignment clause;
%% the inherited Gentzen result remains cited and is not status-upgraded.

\chapter{Arithmetic Proof Theory}
\label{sec:p07-arithmetic-proof-theory}
\chaptermark{Arithmetic Proof Theory}

This chapter records what an ordinal analysis must contain before an ordinal
can be said to calibrate a formal theory.  The inherited source supplied an
honest cited boundary for Gentzen's analysis of arithmetic; it did not supply
the derivation coding, full assignment, or decrease proof.  That status is
preserved here.  The surrounding lookup material uses the ordinal arithmetic
and Cantor normal form developed in
Chapter~\ref{sec:p06-ordinal-arithmetic}.  Primary and modern sources include
\cite{Gentzen1936,Gentzen1938,TroelstraSchwichtenberg2000,
Takeuti1987,Schwichtenberg1977}.

\section{Gentzen's consistency analysis}
\label{sec:p07-gentzen-analysis}

\begin{hbmethod}{def:p07-gentzen-reduction-method}{Gentzen's reduction method}
For the selected sequent presentation of first-order arithmetic, a Gentzen
reduction analysis requires six separately specified objects:
\begin{enumerate}
  \item the exact arithmetic calculus, including its induction inferences;
  \item primitive-recursive codes for its finite derivations;
  \item a primitive-recursive notation system below \(\varepsilon_0\), with
        comparison and fundamental sequences;
  \item a primitive-recursive reduction operation on derivation codes;
  \item an ordinal-notation assignment and a case-by-case theorem that every
        critical reduction strictly lowers it; and
  \item the exact transfinite-induction schema used to exclude an infinite
        reduction, followed by extraction of the absence of a derivation of
        the empty sequent.
\end{enumerate}
The list is necessary because none of its members is determined by the bare
phrase ``ordinal assignment.''  It describes metalinguistic data attached to
one calculus, not additional axioms of arithmetic.  The reduction operation
uses the assignment in the cited termination theorem below; the notation
system alone proves no consistency statement.
\index{Gentzen reduction}
\index{ordinal assignment!derivations}
\end{hbmethod}

\begin{hbtheorem}{thm:p07-gentzen-consistency}{cited=Gentzen1936}{Gentzen's consistency analysis of arithmetic}
Fix the arithmetic sequent-reduction presentation and primitive-recursive
notation system treated in the cited analysis.  Finite derivation codes are
assigned notations below \(\varepsilon_0\) so that every critical reduction
strictly decreases the assigned notation.  In a metatheory accepting
transfinite induction for the predicates and the primitive-recursive ordering
used in that reduction analysis, termination yields the corresponding
arithmetized statement \(\mathsf{Con}(\PA)\).
\end{hbtheorem}

This is the imported result of \cite{Gentzen1936,Gentzen1938}; the exact
modern locator used by the frozen source is
\cite[Chapter~10, especially Sections~10.1--10.4,
pp.~318--337]{TroelstraSchwichtenberg2000}.  The hypothesis is not ``ordinary
induction'' without qualification: it is transfinite induction along the
selected notation ordering for the predicates used in the reduction proof.

\begingroup\small
\begin{longtable}{@{}p{.20\textwidth}p{.36\textwidth}p{.35\textwidth}@{}}
\caption{Components of the cited Gentzen analysis (depth row D-P8-005).}
\label{tab:p07-gentzen-components}\\
\toprule
Component & Exact standing & Evidence retained in this handbook \\
\midrule
\endfirsthead
\toprule
Component & Exact standing & Evidence retained in this handbook \\
\midrule
\endhead
Arithmetic calculus & One named sequent calculus with equality and induction,
not an arbitrary presentation of \(\PA\) & Imported with the cited theorem;
the complete rule surface is not reproduced here \\
Derivation codes & Natural-number codes for finite proof trees, with
primitive-recursive local correctness & Imported; no local coding proof \\
Ordinal notations & Finite codes below \(\varepsilon_0\), with a
primitive-recursive comparison and fundamental-sequence operation & Imported;
Part~6 supplies ordinal meaning, not this coding algorithm \\
Reduction & Primitive-recursive transformation of the selected derivation
codes & Imported; no generic ``proof simplification'' is substituted \\
Assignment and decrease & A map from derivation codes to notation codes and a
strict decrease proof for every critical case & Imported; this is the
load-bearing theorem, not a diagrammatic arrow \\
Induction and extraction & The stated transfinite-induction schema excludes
infinite descent, after which normal-form analysis excludes the empty sequent
& Imported and relative to that metatheory \\
\bottomrule
\end{longtable}
\endgroup

\begin{proofboundary}{def:p07-gentzen-proof-boundary}{Internal boundary of the consistency analysis}
A complete local proof would still have to reproduce the selected arithmetic
rule surface, the derivation coding, the notation grammar and algorithms,
every reduction clause, the complete assignment, every case of the decrease
lemma, the precise transfinite-induction schema, and the final consistency
extraction.  The inherited source explicitly did not do so.  Consequently
Theorem~\ref{thm:p07-gentzen-consistency} is cited, carries no annex item, and
does not claim a proof of \(\mathsf{Con}(\PA)\) formalizable in \(\PA\)
itself.  The second incompleteness theorem blocks that internal conclusion
when \(\PA\) is consistent.
\index{proof boundary!Gentzen analysis}
\end{proofboundary}

\hbstub{N-a-049: the required complete ordinal-assignment table and every decrease case for one selected propositional derivation grammar have not been reproduced source-faithfully; the cited PA analysis and the component ledger above do not fill that must-row remainder.}

\begingroup
\renewcommand{\sectionmark}[1]{\markright{\thesection.\ THEORY PRESENTATIONS}}
\section{Theory presentations used in the calibration table}
\endgroup
\label{sec:p07-calibration-theories}

\begin{hbdefinition}{def:p07-pra}{Primitive recursive arithmetic}
The presentation of \emph{primitive recursive arithmetic} used here, written
\(\mathsf{PRA}\), has a function symbol for every primitive-recursive
numerical function, the corresponding defining equations, classical
equational logic, and quantifier-free induction.  Equivalently presented
free-variable systems may replace this surface only after their equivalence
has been stated.  This definition names a formal first-order theory; it is not
a definition of a philosophical standpoint.  Its ordinal calibration is used
in the first row of Table~\ref{tab:p07-ordinal-calibration}.
\index{primitive recursive arithmetic}
\index{PRA|see{primitive recursive arithmetic}}
\end{hbdefinition}

\begin{hbdefinition}{def:p07-heyting-arithmetic}{Heyting arithmetic}
\emph{Heyting arithmetic}, \(\mathsf{HA}\), has the arithmetic language,
nonlogical axioms, equality axioms, and the Peano-arithmetic induction scheme
fixed in Section~\ref{sec:arithmetic-coding-01}, but uses intuitionistic
first-order logic instead of classical first-order logic.  It is a formal
first-order theory, not a semantic notion of constructive truth.  This
distinction is used by the calibration row here and by every soundness and
derivability statement in Chapter~\ref{sec:p07-functional-interpretations}.
\index{Heyting arithmetic}
\index{HA|see{Heyting arithmetic}}
\end{hbdefinition}

\begin{hbdefinition}{def:p07-isigma-fragments}{The fragments \(\mathsf I\Sigma_n\)}
For each \(n\geq1\), the first-order theory \(\mathsf I\Sigma_n\) is
Robinson arithmetic as fixed in Section~\ref{sec:arithmetic-coding-01} together with
exactly the induction instances whose induction formula lies in the
arithmetical class \(\Sigma^0_n\).  This is a family of formal theories indexed
by \(n\), not one unindexed theory and not the syntactic class
\(\Sigma^0_n\) itself.  The distinction is used in the fragment row of
Table~\ref{tab:p07-ordinal-calibration}.
\index{induction fragment!ISigma-n@\(\mathsf I\Sigma_n\)}
\end{hbdefinition}

\begin{hbdefinition}{def:p07-second-order-calibration-theories}{The second-order calibration theories}
Work in the two-sorted language of second-order arithmetic.  The base theory
\(\mathsf{RCA}_0\) uses classical two-sorted logic, the number-sort axioms of
Robinson arithmetic as fixed in Section~\ref{sec:arithmetic-coding-01}, extensionality
for set variables, induction for \(\Sigma^0_1\) formulas, and comprehension
for \(\Delta^0_1\) formulas.  The theory
\(\mathsf{ACA}_0\) is \(\mathsf{RCA}_0\) plus comprehension for every
arithmetical formula.  The theory \(\mathsf{ATR}_0\) is
\(\mathsf{RCA}_0\) plus the arithmetical-transfinite-recursion scheme along
each coded well-order: every arithmetically defined stage operator has a set
coding its stagewise iteration along that order.

These are axiomatically introduced second-order theories, not alternative
names for \(\PA\) or \(\mathsf{HA}\).  Their set-existence schemes distinguish
them from the first-order theories above, and the distinction is used only in
the two corresponding rows of Table~\ref{tab:p07-ordinal-calibration}
\cite{Simpson2009}.
\index{reverse mathematics!RCA0@\(\mathsf{RCA}_0\)}
\index{reverse mathematics!ACA0@\(\mathsf{ACA}_0\)}
\index{reverse mathematics!ATR0@\(\mathsf{ATR}_0\)}
\end{hbdefinition}

\section{Relative ordinal calibrations}
\label{sec:p07-ordinal-calibration}

\begin{hbdefinition}{def:p07-relative-proof-theoretic-ordinal}{A formula-class-relative proof-theoretic ordinal}
Fix a coding that effectively enumerates recursive binary relations
\(e\mapsto(D_e,\prec_e)\), where each \(D_e\) is a recursive subset of
\(\N\), and fix a class \(\mathcal C\) of formulas.  For a theory \(T\), one
standard first-order calibration is the supremum of
\(\operatorname{otp}(D_e,\prec_e)\) over those \(e\in\N\) for which
\((D_e,\prec_e)\) is externally a well-order and \(T\) proves every instance
of \(\mathcal C\)-transfinite induction along \(\prec_e\).
The index varies across the set whose supremum is taken; fixing one relation
and then reusing it in the set-builder would not define this calibration.  An
order type enters the set only after external well-foundedness of its coded
relation has been established.  This definition names a metalinguistic
calibration relative to the formula class, relation coding, and induction
reading; it is not an intrinsic ordinal object inside \(T\).  The
qualifications are nonredundant: consistency, reflection, provably total
functions, and provable well-foundedness can yield different comparisons.
The table below uses this distinction rather than an unqualified vertical-bar
notation for a theory; vertical bars are reserved for cardinality.
\index{proof-theoretic ordinal!relative calibration}
\index{transfinite induction!ordinal analysis}
\end{hbdefinition}

\begingroup\small
\begin{longtable}{@{}p{.22\textwidth}p{.19\textwidth}p{.50\textwidth}@{}}
\caption{Local ordinal-analysis calibration table (depth row D-P8-006).  Each
row is citation-relative; this is not the Phase-6 master table A-074.}
\label{tab:p07-ordinal-calibration}\\
\toprule
Theory/presentation & Calibrated ordinal & Qualification and source \\
\midrule
\endfirsthead
\toprule
Theory/presentation & Calibrated ordinal & Qualification and source \\
\midrule
\endhead
Modern \(\mathsf{PRA}\), with all primitive-recursive function symbols and
quantifier-free induction & \(\omega^\omega\) & Conventional first-order
induction/provably-total-function calibration; it does not define finitism
\cite{Pohlers2009} \\
\(\mathsf I\Sigma_n\), fixed \(n\geq1\) & \(\omega_{n+1}\), where
\(\omega_0=1\) and \(\omega_{k+1}=\omega^{\omega_k}\) & Standard
fragment calibration for the displayed indexing convention
\cite{Pohlers2009} \\
\(\PA\) & \(\varepsilon_0=\sup_n\omega_n\) & Gentzen
consistency/induction calibration, relative to the notation and metatheory in
Theorem~\ref{thm:p07-gentzen-consistency}
\cite{Gentzen1936,TroelstraSchwichtenberg2000} \\
Heyting arithmetic \(\mathsf{HA}\) & \(\varepsilon_0\) & Standard ordinal
analysis comparison; equality of the displayed calibration does not identify
the classical and intuitionistic proof systems \cite{Pohlers2009} \\
\(\mathsf{ACA}_0\) & \(\varepsilon_0\) & Second-order calibration under the
standard recursive-well-order convention; not a claim that
\(\mathsf{ACA}_0=\PA\) \cite{Simpson2009,Pohlers2009} \\
\(\mathsf{ATR}_0\) and the named autonomous predicative progressions &
\(\Gamma_0\) & The reverse-mathematical and autonomous-progression
calibrations must remain distinguished even when they display the same
ordinal \cite{Feferman1964,Schuette1977,Simpson2009} \\
\bottomrule
\end{longtable}
\endgroup

The table deliberately omits weaker first-order arithmetics, set theories,
and large-cardinal extensions: no number is guessed for them.  Apparatus
A-074 remains a Phase-6 master-table obligation.

\section{Veblen notation and the Feferman--Schuette boundary}
\label{sec:p07-veblen}

\begin{hbdefinition}{def:p07-veblen-hierarchy}{Veblen hierarchy and \(\Gamma_0\)}
An ordinal function is \emph{normal} when it is strictly increasing and is
continuous at every nonzero limit.  Set
\(\operatorname{Veb}_0(\beta)=\omega^\beta\).  For \(\alpha>0\), let
\(\operatorname{Veb}_\alpha\) enumerate in increasing order the common fixed
points of all \(\operatorname{Veb}_\gamma\) with \(\gamma<\alpha\).  The
ordinal \(\Gamma_0\) is
the least positive \(\gamma\) satisfying
\(\operatorname{Veb}_\gamma(0)=\gamma\).

The locally defined symbol \(\operatorname{Veb}\) names this external
hierarchy of ordinal functions; it is distinct from the global formula
metavariables and is needed in the displayed definition of \(\Gamma_0\).
These clauses also define the one ordinal \(\Gamma_0\).  The hierarchy extends
the Cantor-normal-form vocabulary of Part~6 because
fixed points of earlier normal functions must themselves be enumerated.  The
distinction is used in the final row of
Table~\ref{tab:p07-ordinal-calibration}; no ordinal-analysis claim beyond
\(\Gamma_0\) is made here.
\index{Veblen hierarchy}
\index{Gamma-zero@\(\Gamma_0\)}
\end{hbdefinition}

\begin{hbtheorem}{thm:p07-feferman-schuette}{cited=Feferman1964}{Feferman--Schuette predicative calibration}
In the autonomous-progression analyses of Feferman and Schuette,
\(\Gamma_0\) is the bound reached by that formalization of predicative
analysis.
\end{hbtheorem}

The statement is about the cited formal autonomous progressions
\cite{Feferman1964,Schuette1977}; it is not the unqualified philosophical
identity ``predicative mathematics is exactly everything below
\(\Gamma_0\).''  Veblen's original hierarchy is \cite{Veblen1908}; a modern
ordinal-analysis reference is \cite{Pohlers2009}.

\section{Primitive recursive arithmetic and finitism}
\label{sec:p07-pra-finitism}

Skolem's 1923 recursive arithmetic is a historical precursor
\cite{Skolem1923}.  Tait argues that finitistically acceptable numerical
functions and arguments are captured by primitive-recursive methods
\cite{Tait1981}.  The latter is Tait's analysis of finitism, not a theorem
that all philosophical accounts identify finitism with the formal theory
\(\mathsf{PRA}\).

\begin{hbwarning}{def:p07-pra-finitism-boundary}{PRA is not a definition of finitism}
``Finitism'' names a methodological and philosophical class of reasoning;
\(\mathsf{PRA}\) is the exact formal theory in
Definition~\ref{def:p07-pra}.  Skolem's historical program and Tait's later
analysis motivate a comparison, but neither licenses replacing the two names
by an identity sign.  In particular, the unrelated bibliography key
\texttt{Skolem1934} is not evidence for the 1923 arithmetic source.
\index{finitism}
\end{hbwarning}

\hbstub{N-a-052: the could-priority omega-logic entry--the omega-rule, cut elimination for the selected infinitary arithmetic calculus, and its epsilon-zero bound--is not supplied; no infinitary rule is spliced into the finitary Gentzen theorem.}

\paragraph{New-content disposition.}
N-a-049 remains partial under its explicit full-assignment stub.  N-a-050 is
supplied by the exact Veblen-hierarchy definition and the cited, qualified
Feferman--Schuette statement.  N-a-051 is supplied by the fixed modern PRA
surface and the historical/philosophical boundary.  N-a-052 remains at its
required could-row stub.

\paragraph{Source disposition.}
Frozen v1.1 lines 12133--12216 survive as the cited Gentzen statement, the
six-component reading guide, the internal proof boundary, and the warning
against unqualified strength comparisons.  The source's two elementary
fundamental-sequence calculations are absorbed by the Part-6 ordinal entry
and the calibration table; they are not presented as evidence for a
derivation assignment.  The three following source exercises belong to the
handbook's deliberately retired exercise layer and carry no proof status.
